En un ensayo sintético, un tanque bien mezclado de \(100\,\mathrm L\) se lava con agua limpia. El caudal de salida iguala al de entrada y no hay otras pérdidas de soluto. El caudal es constante en cada intervalo de un minuto; en \(t=1\,\mathrm{min}\) se modifica el mando de la bomba, sin retirar ni añadir soluto de forma instantánea.

En cada instante de la tabla se mide consecutivamente una muestra del tanque y un patrón estable de concentración conocida \(5.00\,\mathrm{mg/L}\). El intervalo entre ambas lecturas y los volúmenes de muestra son despreciables. La misma electrónica añade a ambas lecturas un desplazamiento desconocido, que puede cambiar arbitrariamente entre pares. La sensibilidad es uno y cada lectura tiene, además, un error de módulo no superior a \(0.05\,\mathrm{mg/L}\), sin restricciones de signo ni correlación. Las lecturas son:
\[
\begin{array}{c|rrr}
t\ (\mathrm{min})&0&1&2\\\hline
\text{muestra del tanque}\ (\mathrm{mg/L})&16.00&8.60&2.20\\
\text{patrón}\ (\mathrm{mg/L})&5.00&5.60&5.20
\end{array}
\]
Formule el modelo físico y el de observación. Obtenga intervalos garantizados para el caudal antes y después del cambio de mando y decida, por separado, si los datos permiten un único caudal durante los dos minutos y si permiten un desplazamiento electrónico constante. Distinga lo que identifican estas lecturas de lo que queda sin determinar sobre la evolución del desplazamiento. Puede usar calculadora; no se pide un ajuste estadístico.
